

En
effet on peut évidemment supposer @@M0@@ et @@M1@@ affines. Prenons un recouvrement affine @@M2@@ de @@M3@@. Soient @@M4@@~; @@M5@@ un relèvement de @@M6@@ dans @@M7@@. Alors @@M8@@ est égal à la classe du @@M9@@-cocycle @@M10@@~; @@M11@@, à celle de @@M12@@, donc on a
@@M13@@
en tenant compte de l'opération de @@M14@@ sur @@M15@@, donc sur @@M16@@.

D'après ii), on a @@M17@@ et donc @@M18@@ pour
@@M19@@ et
@@M20@@. Or par (v) @@M21@@ est surjectif~; par suite @@M22@@ pour
@@M23@@, @@M24@@ et
@@M25@@.

D'après i),
@@M26@@. Par suite @@M27@@ où @@M28@@ d'où @@M29@@, donc @@M30@@.

Finalement soit @@M31@@. Soit @@M32@@ où
@@M33@@; on a alors @@M34@@, d'où @@M35@@.

En vertu de 3.2, iii) et iv), on a une suite exacte fonctorielle en @@M36@@~:
@@M37@@
qui par passage aux faisceaux f.p.p.f. associés donne une suite exacte
@@M38@@

Si @@M39@@, le monomorphisme @@M40@@ s'obtient par changement de base à partir de la section unité @@M41@@~; comme @@M42@@ est représentable et séparé, il résulte que @@M43@@ est représentable par une
